% ============================================================
% Section 91: 自由电子气的压强与体积弹性模量
% ============================================================
\section{固体理论：自由电子气的压强与体积弹性模量}
\label{sec:free-electron-pressure}

\begin{modelbox}[题目：自由电子气压强、弹性模量的推导与数值估算]
在绝对零度下：
\begin{enumerate}
    \item 证明金属自由电子气的压强 $p=\dfrac{2}{3}\dfrac{U_0}{V}$，其中 $U_0$ 为电子气总能量，$V$ 为金属体积；
    \item 证明电子气的体积弹性模量 $K=-V\dfrac{\partial p}{\partial V}=\dfrac{10}{9}\dfrac{U_0}{V}$；
    \item 钾金属电子浓度为 $n=1.40\times10^{22}\,\text{cm}^{-3}$，计算 $p$ 和 $K$。
\end{enumerate}
\end{modelbox}

\begin{tipbox}[考点快速分析]
\begin{itemize}
    \item \textbf{热力学关系}：$p=-(\partial U_0/\partial V)_N$（$T=0$ 时 $dS=0$）
    \item \textbf{费米能的体积依赖}：$E_F\propto V^{-2/3}$，$U_0\propto V^{-2/3}$
    \item \textbf{弹性模量}：$K=-V\partial p/\partial V$，反映抗压缩性
    \item \textbf{数量级}：电子简并压达数万大气压
\end{itemize}
\end{tipbox}

\begin{derivationbox}[标准解题过程]
\textbf{Step 1: 压强的证明}

$T=0$ 时，$dS=0$，热力学关系给出
\begin{equation}
p=-\Bigl(\frac{\partial U_0}{\partial V}\Bigr)_N.
\end{equation}

自由电子气中，$T=0$ 时每个电子平均能量
\begin{equation}
\bar E=\frac{U_0}{N}=\frac{3}{5}E_F^0,
\end{equation}
其中 $E_F^0=\dfrac{\hbar^2}{2m}(3\pi^2N/V)^{2/3}$。总能量
\begin{equation}
U_0=\frac{3}{5}N\frac{\hbar^2}{2m}\Bigl(\frac{3\pi^2N}{V}\Bigr)^{2/3}=CV^{-2/3},
\end{equation}
$C$ 与 $V$ 无关。

求导得
\begin{equation}
p=-C\Bigl(-\frac{2}{3}\Bigr)V^{-5/3}=\frac{2}{3}CV^{-2/3}
=\frac{2}{3}\frac{U_0}{V}.
\end{equation}
故
\begin{equation}
\boxed{p=\frac{2}{3}\frac{U_0}{V}=\frac{2}{5}nE_F^0.}
\end{equation}

\textbf{Step 2: 体积弹性模量的证明}

体积弹性模量
\begin{equation}
K=-V\frac{\partial p}{\partial V}.
\end{equation}

由 $U_0\propto V^{-2/3}$ 知 $p=\dfrac{2}{3}\dfrac{U_0}{V}\propto V^{-5/3}$，故
\begin{equation}
\frac{\partial p}{\partial V}=-\frac{5}{3}\frac{p}{V}.
\end{equation}

代入得
\begin{equation}
K=-V\Bigl(-\frac{5}{3}\frac{p}{V}\Bigr)=\frac{5}{3}p
=\frac{5}{3}\cdot\frac{2}{3}\frac{U_0}{V}
=\boxed{\frac{10}{9}\frac{U_0}{V}=\frac{2}{3}nE_F^0.}
\end{equation}

\textbf{Step 3: 钾金属的 $p$ 和 $K$}

已知 $n=1.40\times10^{28}\,\text{m}^{-3}$。先算费米能
\begin{equation}
E_F^0=\frac{\hbar^2}{2m}(3\pi^2n)^{2/3}
\approx\frac{(1.055\times10^{-34})^2}{2\times9.11\times10^{-31}}\times(4.145\times10^{29})^{2/3}
\approx3.40\times10^{-19}\,\text{J}\approx2.12\,\text{eV}.
\end{equation}

压强
\begin{equation}
\boxed{p=\frac{2}{5}nE_F^0=\frac{2}{5}\times1.40\times10^{28}\times3.40\times10^{-19}
\approx1.90\times10^9\,\text{Pa}\approx1.90\times10^4\,\text{atm}.}
\end{equation}

体积弹性模量
\begin{equation}
\boxed{K=\frac{5}{3}p\approx\frac{5}{3}\times1.90\times10^9
\approx3.17\times10^9\,\text{Pa}\approx3.17\times10^4\,\text{atm}.}
\end{equation}
\end{derivationbox}

\begin{formulabox}[答案汇总]
\begin{center}
\begin{tabular}{lll}
\toprule
\textbf{问题} & \textbf{表达式} & \textbf{数值（钾）} \\
\midrule
(1) 压强 & $p=\dfrac{2}{3}\dfrac{U_0}{V}=\dfrac{2}{5}nE_F^0$ & $1.90\times10^9\,\text{Pa}$ \\
(2) 弹性模量 & $K=\dfrac{10}{9}\dfrac{U_0}{V}=\dfrac{2}{3}nE_F^0$ & $3.17\times10^9\,\text{Pa}$ \\
\bottomrule
\end{tabular}
\end{center}
\end{formulabox}

\begin{insightbox}[物理图像]
\begin{itemize}
    \item 自由电子气具有巨大的\textbf{简并压}（数万大气压），源于泡利不相容原理导致的基态零点动能。即使在绝对零度，电子气仍对外施加巨大压力，这是金属具有有限平衡体积的重要原因之一。
    \item 体积弹性模量 $K$ 是金属\textbf{抗压缩性}的量度。碱金属的 $K$ 主要由自由电子气贡献；对于过渡金属，$d$ 电子的库仑关联也会显著贡献 $K$。
    \item 本题将金属的宏观力学性质（压强、弹性模量）与微观电子浓度联系起来，是理解金属状态方程的基础。实际金属的总压强还需包含离子实的排斥贡献，二者平衡决定晶格常数。
\end{itemize}
\end{insightbox}
